\subsubsection{Three-by-six bands with five outside elements}

\begin{lemma}\label{B:five-two-by-five-pairs}
In the two exceptional algebras of
Proposition~\ref{B:five-nonlocal-exceptions}, suppose
\[
 \dim V=2,\quad\dim W=5,\qquad
 (\dim\operatorname{Tor}_0,
   \dim\operatorname{Tor}_1,
   \dim\operatorname{Tor}_2)=(2,0,1).
\]
For the zero-composite line algebra the only pair is
\[
 V=S_1\oplus S_3,\qquad
 W=T_1\oplus e_1A\oplus e_2A.
\]
Every balanced form for this pair factors through the same
one-dimensional quotient $S_1$ of $V$.
For the Nakayama algebra the only pair is
\[
 V=L=Ae_1/\mathbb Qc,\qquad
 W=T_2\oplus(e_2A)^{\oplus2}.
\]
Both modules of this latter pair are faithful over
$A/I$, where $I=\operatorname{span}(a,c)$, and
$A/I$ is the path algebra of one arrow.
\end{lemma}
\begin{proof}
Use the decompositions and exact Tor table in
Lemma~\ref{B:five-exceptional-targets}.
For the line algebra the same degree-one and degree-two
conditions force $V=S_1\oplus S_3$ and
\[
 W=T_1\oplus T_3^{\oplus s}\oplus(e_1A)^{\oplus p}
                   \oplus(e_2A)^{\oplus q},\qquad
 s+2p+2q=4.
\]
The tensor dimension is $1+s+p=2$, whose unique solution is
$(s,p,q)=(0,1,1)$.  Both nonzero tensor summands use only
the $S_1$ coordinate of $V$, proving the common factorization.

For the Nakayama algebra the degree-one and degree-two
conditions force $V=L$ and
$W=T_2\oplus(e_1A)^{\oplus p}\oplus(e_2A)^{\oplus q}$.
Now $3p+2q=4$, so $(p,q)=(0,2)$; its tensor dimension is two.
The actions of $1,e_1,b$ on $L$ are linearly independent, and
the same is true on the right module $e_2A$.  Their kernels
as $A$-representations are therefore exactly
$I=\operatorname{span}(a,c)$.  Quotienting by this ideal
leaves two vertex idempotents and one arrow, as claimed.
\end{proof}

\begin{lemma}[The actual Nakayama multiplicative bases]
\label{B:five-Nakayama-common-annihilator}
Suppose a contracted finite-monoid algebra of dimension five
is isomorphic to the Nakayama algebra in
Lemma~\ref{B:five-two-by-five-pairs}.  In its distinguished
monoid basis, the ideal $I=\operatorname{span}(a,c)$ is spanned
by basis vectors lying in $I$ and differences of two basis
vectors that have the same image in $A/I$.
\end{lemma}
\begin{proof}
This is a finite multiplicative-basis assertion, checked against
the complete independently generated order-six census, including
all reducible monoids and all monoids with zero.
The checker \path{code/d_nakayama_five_annihilator.py} examines
all $2325$ distinct literal table/opposite choices supplied by
the $1373$ isomorphism-and-opposite classes in
\path{code/outputs/c_monoids_6.tsv}.

For clarity, the algebra-recognition and ideal test are exact.
After contracting the unique zero, enumerate all rational
characters, whose values on a finite monoid belong to
$\{-1,0,1\}$.  Require exactly two independent characters;
their common kernel $J$ has dimension three.  Require
$\dim J^2=1$ and $J^3=0$, so this kernel is exactly the
radical.  Lift the two orthogonal idempotents and order them
so that the radical corner dimensions are
\[
 (\dim e_1Je_1,\dim e_1Je_2,
   \dim e_2Je_1,\dim e_2Je_2)=(1,1,1,0).
\]
The additional tests
$(e_1Je_2)(e_2Je_1)\ne0$ and
$(e_2Je_1)(e_1Je_2)=0$ identify precisely the stated
Nakayama algebra.  This finds fourteen oriented bases, from
seven source census classes.

In each case the checker forms the intrinsic ideal
$I=e_1Je_2+J^2$, computes the images of every distinguished
basis vector in the three-dimensional quotient, and spans
the zero-image vectors and equal-image differences.  Exact
rational row reduction shows that this span equals $I$
in all fourteen cases.  The full tables, idempotent-independent
ideal bases, quotient coordinates, and spanning relations are
recorded in
\path{code/outputs/d_nakayama_five_annihilator.json}.
There is no assumption that an abstract algebra isomorphism
preserves a particular displayed monoid basis.
An independent implementation,
\path{code/b_verify_nakayama_five_bases.py}, repeats the entire
classification using rational \texttt{Fraction} arithmetic and a
different idempotent lift.  It finds the same fourteen oriented
bases and verifies every ideal-spanning equality; its results are
in \path{code/outputs/b_nakayama_five_bases_verified.json}.
\end{proof}

\begin{proposition}\label{B:three-by-six-five-outside}
A finite monoid with a $3\times6$ or $6\times3$ rectangular-band
minimal ideal and exactly five elements outside that ideal
cannot have classifying space homotopy equivalent to
$S^2\times S^2$.
\end{proposition}
\begin{proof}
Take the opposite if necessary so that the augmentation modules
have dimensions $\dim V=2$, $\dim W=5$, and put
$A=\mathbb QM/\mathbb QK$.  The band Tor formula gives the
Tor dimensions $2,0,1,0,\ldots$.  A local quotient is excluded
by Corollary~\ref{B:local-three-by-six-product}.  Otherwise
Proposition~\ref{B:five-nonlocal-exceptions} and eventual Tor
vanishing leave only the two pairs in
Lemma~\ref{B:five-two-by-five-pairs}.

For the line pair every balanced form has a common
one-dimensional factor.  Lemma~\ref{B:rank-one-square} and
polarization make all degree-two products zero.

For the Nakayama pair, let $H$ be the faithful transformation-pair
image of $M$.  We claim that the induced map of contracted
rational algebras kills $I$.  Each basis vector in $I$ has zero
row and column augmentation operators, so its two set maps
are constant and its profile belongs to the minimal band.
For two basis vectors equal modulo $I$, their augmentation
operators agree on both sides.  Since both $V$ and $W$ are
faithful over $A/I$, the common operators are either both
nonzero or both zero.  In the former case each pair of set
maps is equal: two finite set maps with the same nonzero
augmentation operator must coincide.  To see this directly,
equality of the operators makes
$[f(x)]-[g(x)]$ independent of $x$; a nonzero such difference
forces both maps to be constant.  In the zero case both
profiles already lie in the band.  Thus the image kills every
spanning relation in Lemma~\ref{B:five-Nakayama-common-annihilator},
and hence kills $I$.

The contracted algebra of $H$ is consequently a quotient of
the path algebra of one arrow.  Every quotient of this specific
algebra is hereditary: it either retains the single arrow or
is a product of copies of $\mathbb Q$.  The band Tor formula
gives $H_4(BH;\mathbb Q)=0$.  The natural map
$H_2(BM;\mathbb Q)\to H_2(BH;\mathbb Q)$ is an isomorphism,
because the balanced action presentations are identical.
All degree-two classes of $M$ therefore pull back from $H$,
and all their products vanish.  This again contradicts the
cohomology ring of $S^2\times S^2$.
\end{proof}

\begin{remark}\label{B:Nakayama-two-by-five-rank-correction}
The Nakayama balanced forms themselves have rank two.  On one
$e_2A$ summand, in the bases $(e_2,b)$ and $(\bar e_1,\bar b)$,
the multiplication pairing is
$\left(\begin{smallmatrix}0&1\\1&0\end{smallmatrix}\right)$.
The proof uses the actual multiplicative-basis quotient and
does not infer rank-one factorization from
$e_2A\otimes_A L\cong e_2L$.
As a separate construction check,
\path{code/d_product23_nakayama.py} enumerates all $732$
affine lifts for the standard zero basis: four row-action
types and $183$ column-action types.  Every resulting
twenty-three-element monoid has the required finite additive
homology, but all $2\times2$ cup matrices are zero.  This
last family calculation is not used in the exhaustive
multiplicative-basis argument above.
\end{remark}
